\documentclass[11pt,a4paper]{article}

\usepackage[utf8]{inputenc}
\usepackage[T1]{fontenc}
\usepackage{lmodern}
\usepackage{amsmath,amssymb,amsthm,mathtools}
\usepackage{geometry}
\geometry{margin=2.8cm}
\usepackage{hyperref}
\usepackage{enumitem}
\usepackage{xcolor}
\usepackage{tcolorbox}
\tcbuselibrary{breakable,skins}
\usepackage{microtype}

\newtheorem{theorem}{Theorem}[section]
\newtheorem{lemma}[theorem]{Lemma}
\newtheorem{proposition}[theorem]{Proposition}
\newtheorem{corollary}[theorem]{Corollary}
\newtheorem{conjecture}[theorem]{Conjecture}
\theoremstyle{definition}
\newtheorem{definition}[theorem]{Definition}
\newtheorem{remark}[theorem]{Remark}
\newtheorem{example}[theorem]{Example}

\newcommand{\N}{\mathbb{N}}
\newcommand{\Z}{\mathbb{Z}}
\newcommand{\Q}{\mathbb{Q}}
\newcommand{\R}{\mathbb{R}}
\newcommand{\C}{\mathbb{C}}
\newcommand{\F}{\mathbb{F}}
\newcommand{\K}{\mathbb{K}}
\newcommand{\sub}{\mathrm{sub}}
\newcommand{\Hecke}{\mathbb{T}}
\newcommand{\fp}{\mathrm{fund}}
\newcommand{\Vmin}{V_{\min}}
\newcommand{\Tp}{T_p}
\newcommand{\Tn}{T_n}
\newcommand{\Wit}{\mathrm{Wit}}
\newcommand{\Cphi}{\mathcal{C}_{\!\varphi}}

\hypersetup{
  colorlinks=true,
  linkcolor=blue!60!black,
  citecolor=blue!60!black,
  urlcolor=blue!60!black,
}

\title{\textbf{The Fractal Tiling Theorem}\\
\large Reducing RH to a Finite Hecke Verification via Substrate Multiplicativity}
\author{Lee Smart \\ \small Institute of Vibrational Field Dynamics}
\date{2026-05-30 \quad v1.0.0-rc1}

\begin{document}
\maketitle

\begin{center}
\fbox{\parbox{0.92\textwidth}{\centering
\textsc{Pre-peer-review open research preprint.} The substrate-completeness step is
conditional on finite cuspidal Hilbert modular data over $\Q(\sqrt{5})$ that has
not yet been constructed unconditionally. \textbf{No unconditional proof of RH
is claimed.} The contribution is a clean conditional reduction: if the substrate
provides admissible Hecke data at finitely many primes, the Riemann Hypothesis
follows by tiling.
}}
\end{center}

\begin{abstract}
We isolate three independent symmetry structures on the $V_{600}$ closure
substrate: the geometric tiling by $W(H_4)$, the cascade scale-tiling
$\mathrm{Pentagon}\!\to\!\mathrm{Icos}\!\to\!V_{600}\!\to\!E_8\!\to\!\cdots$, and
the Hecke multiplicative tiling generated by prime operators $\{T_p\}_{p\,\mathrm{prime}}$.
The first two extend statements across the substrate's interior; only the third
extends statements to the analytical content of the L-function.
We prove the \emph{Fractal Tiling Theorem}: if the substrate's witness
eigenspace $V_{\min}\subset V_{600}$ supports Hecke operators $\{T_p\}_{p\le P}$
that satisfy the seven admissibility classes of the substrate translation
engine, then by Hecke multiplicativity these tile to all $T_n$ and the tiled
Dirichlet series is a genuine cuspidal automorphic L-function. We then
\emph{carry out} the construction at level $\mathfrak{p}_{31}$: using the
icosian ring as the maximal order of $(-1,-1/\Q(\sqrt5))$, the substrate's own
Brandt/Hecke action reproduces the genuine norm-$31$ Hilbert newform exactly
(out-of-sample, no fitting, $11$ prime ideals), and the verifiable admissibility
classes (multiplicativity, Ramanujan, Sato-Tate) pass on this real data.
\textbf{We also correct an overclaim:} admissibility plus the Weil explicit
formula certify automorphy but do \emph{not} place the zeros on the critical
line --- RH for the resulting L-function is equivalent to GRH and remains open.
The honest contribution is therefore a verified, non-circular substrate$\leftrightarrow$Hecke
identification, not a proof of RH.
\end{abstract}

\tableofcontents

\section{Introduction}

\subsection{The shortcut question}
The substrate framework developed in prior bundles establishes:
\begin{itemize}
\item the closure operator $\Cphi = (12+\varphi^{-2})I - A_1$ on $V_{600}$
      \cite{IcosianTriad};
\item the substrate L-function identity $L_\sub(s) = \zeta_K(s)\,\zeta_K(s-1)\,C_2(s)$
      with $K=\Q(\sqrt{5})$ \cite{CriticalLinePullback};
\item the Witness Invariant Theorem $W = \Vmin$ \cite{ObserverAttractor};
\item the spectral-resonance reformulation $\mathrm{RH}\Leftrightarrow T_W$
      exists with spectrum $\{\gamma_n\}$ \cite{RHWitnessResonance};
\item the gradient-bridge cascade-coverage analysis: $V_{600}$ alone provides
      $\sim\!65\%$ coverage of the off-axis topology, with the infinite-cascade
      limit required for $100\%$ \cite{GradientBridge}.
\end{itemize}
Naively, this leaves RH a difficult open problem because no finite cascade level
reaches the analytical content of off-axis zeros directly. This paper isolates a
\emph{shortcut} that does not require finite cascade coverage at the level of
the strip topology. The shortcut is multiplicative tiling of Hecke data,
combined with the Weil explicit formula.

\subsection{The three tilings}
Three structures on the substrate act as fundamental-domain-with-extension
mechanisms:

\begin{enumerate}[label=(T\arabic*)]
\item \textbf{Geometric tiling.} $W(H_4)$ acts on $V_{600}$ with order
      $14{,}400=120^2$, single orbit on the 120 vertices, vertex stabiliser of
      order $120$. Any $W(H_4)$-equivariant statement on one vertex extends to
      all $120$.

\item \textbf{Cascade scale-tiling.} The cascade
      $\mathrm{Pentagon}\to\mathrm{Icos}\to V_{600}\to E_8\to E_8\!\times\!E_8\to\cdots$
      is a sequence of McKay-related binary lifts. Algebraic patterns at one
      level propagate to all higher levels by lifting/projection.

\item \textbf{Hecke tiling.} The Hecke algebra $\Hecke$ at fixed level is
      multiplicatively generated by operators indexed by rational primes:
      $T_m T_n = T_{mn}$ for $\gcd(m,n)=1$, with a polynomial recursion at
      prime powers. Therefore $\{T_n\}_{n\in\N}$ is determined by
      $\{T_p\}_{p\,\mathrm{prime}}$.
\end{enumerate}

The first two tilings are \emph{intra-substrate}: they extend statements within
$V_{600}$ and the cascade. The third tiling is the one that extends substrate
data to the L-function: through it, finite-prime substrate Hecke data forces
constraints on all Riemann zeros via the Weil explicit formula.

\subsection{Result}
\begin{tcolorbox}[colback=blue!2,colframe=blue!50!black,arc=2mm,
    title={\textbf{Fractal Tiling Theorem (informal)}},fonttitle=\bfseries]
If the substrate's witness eigenspace $\Vmin$ supports a finite-prime Hecke
system $\{T_p\}_{p\le P}$ that simultaneously satisfies the seven admissibility
classes of the substrate translation engine (Ramanujan, multiplicativity,
functional equation, Galois equivariance, real positivity, explicit-formula
convergence, and Sato-Tate distribution), then the tiled $L_\sub$ is a genuine
cuspidal automorphic L-function. \emph{(This is what the admissibility classes
certify. It does \textbf{not} yield RH for $L_\sub$: see
Remark~\ref{rem:no-rh}. RH for the resulting L-function is equivalent to GRH
and remains open.)}
\end{tcolorbox}

The substrate has already established $W=\Vmin$ unconditionally. What remains
is the construction of $\{T_p\}_{p\le P}$ \emph{as Hecke operators on a genuine
cuspidal automorphic representation} of $\mathrm{GL}_2(\Q(\sqrt{5}))$ whose
eigenvalues coincide with the substrate's spectral data on the 26-dimensional
block of $V_{600}$. This is the precise finite mathematical task remaining.

\section{The Hecke fundamental domain}

\subsection{Setup}
Fix the substrate $V_{600}$ with closure operator $\Cphi$. The witness space is
the one-dimensional eigenspace
$$
\Vmin := \ker(\Cphi - \varphi^{-2} I)\quad\text{(when restricted to the
$\tau$-fixed block)}.
$$
Let $K=\Q(\sqrt{5})$, $\mathfrak{O}_K = \Z[\varphi]$, and let
$\mathfrak{N}\subset\mathfrak{O}_K$ be the level ideal forced by the substrate's
icosian discriminant. The genuine level is a \emph{prime ideal}
$\mathfrak{N}=\mathfrak{p}_{31}\subset\mathfrak{O}_K$ of norm $31$ (the prime
$31\equiv1\pmod 5$ splits in $K$, so such an ideal exists); this is the level
at which the first cuspidal Hilbert newform over $\Q(\sqrt 5)$ appears
\cite{Dembele}, matching the proxy computation's $N=31$.
\emph{(Earlier drafts wrote $\mathfrak{N}=(2)\cdot(\varphi-2)$; this is an
error, since $\varphi-2=-\varphi^{-2}$ is a unit, so that product collapses to
$(2)$, which is inert of norm $4$ and supports no cusp form.)}

For each prime ideal $\mathfrak{p}\subset\mathfrak{O}_K$ coprime to $\mathfrak{N}$
the substrate's spectral data on the $26$-dimensional block carries a candidate
Hecke action $\Tp^{\mathrm{sub}}$ defined by the closure-flow's eigenvalue at
the $\mathfrak{p}$-graded subspace.

\subsection{The fundamental domain}
\begin{definition}[Hecke fundamental domain]\label{def:fund}
The \emph{Hecke fundamental domain} for the substrate at level $\mathfrak{N}$
and bound $P$ is the finite set
$$
\mathcal{F}_P :=
\{\Tp^{\mathrm{sub}} : \mathfrak{p}\subset\mathfrak{O}_K,\ \mathfrak{p}\nmid\mathfrak{N},\ N(\mathfrak{p})\le P\}.
$$
\end{definition}

This is a finite collection of self-adjoint operators on the substrate's
26-dim block. Its size is $\pi_K(P)$, the prime ideal counting function over
$K$. For $K=\Q(\sqrt{5})$ and $P=200$ one has $\pi_K(P)=46$ ideals: the
fundamental domain has only $46$ generators.

\subsection{Multiplicativity tiling}
\begin{theorem}[Multiplicativity tiling]\label{thm:mult-tiling}
Let $\{\Tp^{\mathrm{sub}}\}_{N(\mathfrak{p})\le P}$ be a collection of
operators on a finite-dimensional Hilbert space, all simultaneously
diagonalisable, with common eigenform $\psi$ and eigenvalues $a_\mathfrak{p}$.
Suppose
\begin{enumerate}[label=(\roman*)]
\item the $a_\mathfrak{p}$ satisfy $|a_\mathfrak{p}|\le 2\sqrt{N(\mathfrak{p})}$ (Ramanujan);
\item the Hecke multiplicativity recursion holds:
$$
a_{\mathfrak{p}^{k+1}}
= a_\mathfrak{p}\,a_{\mathfrak{p}^k} - N(\mathfrak{p})\,a_{\mathfrak{p}^{k-1}},
$$
\item for coprime ideals $\mathfrak{a}\perp\mathfrak{b}$: $a_{\mathfrak{a}\mathfrak{b}} = a_\mathfrak{a}\,a_\mathfrak{b}$.
\end{enumerate}
Then $\{a_\mathfrak{n}\}_{\mathfrak{n}\subset\mathfrak{O}_K}$ is uniquely
determined for all ideals $\mathfrak{n}$ coprime to $\mathfrak{N}$ from the
finite data $\mathcal{F}_P$ by induction on $N(\mathfrak{n})$. Moreover the
Dirichlet series
$$
L(s,\psi) := \sum_{\mathfrak{n}\perp\mathfrak{N}} \frac{a_\mathfrak{n}}{N(\mathfrak{n})^s}
= \prod_{\mathfrak{p}\nmid\mathfrak{N}}
   \left(1 - a_\mathfrak{p}\,N(\mathfrak{p})^{-s} + N(\mathfrak{p})^{1-2s}\right)^{-1}
$$
converges absolutely for $\Re s > 3/2$ and the Euler product is unconditional.
\end{theorem}

\begin{proof}
The induction is the standard Hecke recursion: for $\mathfrak{n} = \prod_i
\mathfrak{p}_i^{k_i}$ with $\mathfrak{p}_i$ distinct, multiplicativity gives
$a_\mathfrak{n} = \prod_i a_{\mathfrak{p}_i^{k_i}}$, and each $a_{\mathfrak{p}_i^{k_i}}$
is computed from $a_{\mathfrak{p}_i}$ by the recursion in (ii). Ramanujan (i)
bounds $|a_\mathfrak{n}|\le d(\mathfrak{n})\sqrt{N(\mathfrak{n})}$ where
$d(\mathfrak{n})$ is the ideal divisor function, giving the absolute convergence
of the Dirichlet series for $\Re s > 3/2$. The Euler product follows from
multiplicativity by the standard formal identity.
\end{proof}

\begin{remark}\label{rem:tiling}
Theorem \ref{thm:mult-tiling} is the substrate's \textbf{fractal tiling}: the
finite data $\mathcal{F}_P$ generates the entire infinite Dirichlet series
$L(s,\psi)$. Doubling $P$ doesn't double the work because the new primes don't
introduce qualitatively new data --- they're constrained by multiplicativity.
The fundamental domain is finite. The L-function is the tiled extension.
\end{remark}

\section{The explicit formula bridge}

\subsection{Statement}
The Weil explicit formula \cite{Weil}, in the form adapted to Hilbert modular
L-functions over $K=\Q(\sqrt{5})$, asserts: for any test function $h$ in the
admissible class (Paley-Wiener with rapid decay) and $\hat h$ its Fourier
transform,
$$
\sum_{\rho} h(\gamma_\rho) =
\underbrace{\hat h(0) \log\!\left(\tfrac{|d_K|\,N(\mathfrak{N})}{\pi^4}\right)
 + 2 \int_{-\infty}^{\infty} h(r) \,\Psi(r)\,dr}_{\text{archimedean side}}
 - 2\sum_{\mathfrak{p}\nmid\mathfrak{N}}
       \sum_{k\ge 1}
       \frac{a_{\mathfrak{p}^k}\,\log N(\mathfrak{p})}{N(\mathfrak{p})^{k/2}}
       \,\hat h(k\log N(\mathfrak{p}))
$$
where $\rho = 1/2 + i\gamma_\rho$ ranges over non-trivial zeros of $L(s,\psi)$
and $\Psi$ is the standard archimedean density.

\subsection{The bridge}
\begin{lemma}[Explicit formula bridge]\label{lem:bridge}
Under the hypotheses of Theorem \ref{thm:mult-tiling}, the right-hand side of
the explicit formula is fully determined by the finite-prime data
$\mathcal{F}_P$ together with the multiplicativity tiling. Specifically, the
infinite sum over primes can be evaluated to arbitrary precision $\epsilon$ by
truncating to $N(\mathfrak{p})\le P_0(h,\epsilon)$ for an explicit bound
$P_0(h,\epsilon)$ depending only on the rapid decay of $\hat h$ and on
$\epsilon$.
\end{lemma}

\begin{proof}
Ramanujan gives $|a_{\mathfrak{p}^k}|\le (k+1)\sqrt{N(\mathfrak{p})^k}$ so that
the inner sum over $k$ is controlled by a geometric series:
$$
\sum_{k\ge 1}\frac{|a_{\mathfrak{p}^k}|}{N(\mathfrak{p})^{k/2}}|\hat h(k\log N(\mathfrak{p}))|
\le \sum_{k\ge 1}(k+1)|\hat h(k\log N(\mathfrak{p}))|.
$$
Rapid decay of $\hat h$ then bounds the tail $N(\mathfrak{p})>P_0$ by
$\epsilon$. The remaining finite sum involves only $a_\mathfrak{p}$ for
$N(\mathfrak{p})\le P_0$, which by Theorem \ref{thm:mult-tiling} is determined
by the fundamental domain $\mathcal{F}_{P_0}$.
\end{proof}

\begin{corollary}\label{cor:zeros-from-tiling}
The locations of the zeros $\gamma_n$ of $L(s,\psi)$, sampled by the test
function $h$, are determined to precision $\epsilon$ by $\mathcal{F}_{P_0(h,\epsilon)}$.
\end{corollary}

\section{The seven admissibility classes}\label{sec:admissibility}

The substrate translation engine v0.6 \cite{TranslationEngine} verifies
candidate finite Hecke data against seven admissibility classes. We list them
here in the form they appear in the engine's verdict layer; each is a
necessary condition on \emph{genuine} cuspidal Hecke data.

\begin{enumerate}[label=(A\arabic*)]
\item \textbf{(SM) Substrate multiplicativity.} The Hecke recursion at prime
      powers holds exactly.
\item \textbf{(FE) Functional equation.} The completed L-function satisfies
      $\Lambda(s,\psi) = \pm\Lambda(2-s,\psi)$.
\item \textbf{(DN) Dirichlet normalisation.} The leading coefficient
      $a_{(1)}=1$.
\item \textbf{(GUE) Galois universal equivariance.} The Galois action
      $\tau\colon\varphi\mapsto 1-\varphi$ acts compatibly on eigenvalues.
\item \textbf{(EF) Explicit formula convergence.} The Weil explicit formula
      computed from $\mathcal{F}_P$ converges as $P\to\infty$ to the
      archimedean side, with the tail error bounded by Lemma \ref{lem:bridge}.
\item \textbf{(CB) Cuspidal bound.} $|a_\mathfrak{p}| < 2\sqrt{N(\mathfrak{p})}$
      strictly (Ramanujan-bound, cuspidal, not Eisenstein).
\item \textbf{(ST) Sato-Tate distribution.} For varying $\mathfrak{p}$ in
      $\mathcal{F}_P$, the normalised eigenvalues
      $a_\mathfrak{p}/(2\sqrt{N(\mathfrak{p})})\in[-1,1]$ are statistically
      consistent with the semicircle density $\frac{2}{\pi}\sqrt{1-x^2}$.
\end{enumerate}

\begin{remark}
Of the seven, (CB) is the most consequential: it distinguishes cuspidal forms
(needed for RH) from Eisenstein forms (which sit at trivial loci). In the
engine's level-31 BrandtModule(31,1) computation over $\Q$, raw icosian theta
eigenvalues failed (CB) until the Eisenstein projection $a_p\mapsto a_p - (p+1)$
was applied. Over the genuine icosian setting at level $\mathfrak{N}\subset\Q(\sqrt{5})$,
the projection's analogue is the closure-flow's removal of the constant mode.
\end{remark}

\section{The Fractal Tiling Theorem}

\subsection{Statement}
\begin{theorem}[Fractal Tiling Theorem]\label{thm:fractal-tiling}
Let $\Vmin\subset V_{600}$ be the witness eigenspace of the closure operator
$\Cphi$, restricted to the $\tau$-fixed block. Suppose there exists a finite
set of self-adjoint operators
$$
\mathcal{F}_P = \{T_\mathfrak{p}^{\mathrm{sub}} : \mathfrak{p}\subset\mathfrak{O}_K,\ \mathfrak{p}\nmid\mathfrak{N},\ N(\mathfrak{p})\le P\}
$$
acting on the $26$-dimensional block of $V_{600}$, with a common eigenform
$\psi\in\Vmin$, eigenvalues $a_\mathfrak{p}$, satisfying simultaneously the
admissibility classes (A1)--(A7) of Section \ref{sec:admissibility}.

Then the Dirichlet series $L(s,\psi)$ tiled from $\mathcal{F}_P$ by Theorem
\ref{thm:mult-tiling}:
\begin{enumerate}
\item agrees with $L_\sub(s) = \zeta_K(s)\,\zeta_K(s-1)\,C_2(s)$ up to the
      finite primes above $\mathfrak{N}$;
\item is a genuine cuspidal Hilbert modular L-function on
      $\mathrm{GL}_2(\Q(\sqrt{5}))$.
\end{enumerate}
\end{theorem}

\begin{remark}[Important scope correction]\label{rem:no-rh}
An earlier version of this theorem asserted a third conclusion --- that the
tiled L-function \emph{satisfies the Riemann Hypothesis}. \textbf{That
conclusion is withdrawn.} As shown by the explicit construction of
\S\ref{sec:status} and the companion file \texttt{route\_b/admissibility.py},
the admissibility classes (A1)--(A7) certify that $L(\psi,s)$ is a genuine
\emph{automorphic} L-function, but they do \emph{not} imply that its zeros lie
on the critical line. RH for $L(\psi,s)$ is precisely GRH for the corresponding
object (here, an elliptic curve over $\Q(\sqrt5)$), which is open. The Weil
explicit formula is a valid identity relating zeros to primes, but it is a tool
equivalent to information about the zeros --- it does not force them onto the
line. The reduction below is therefore a reduction of the \emph{L-function}
to finite data, not a reduction of \emph{RH} to finite data.
\end{remark}

\subsection{Proof outline}
\begin{proof}
\textbf{(1) Identification with $L_\sub$.}
By (A1)--(A2)+(A4), $\mathcal{F}_P$ generates a Dirichlet series with the
correct functional equation and Galois-equivariance, matching the symbolic
$L_\sub$ on its prime factorisation away from $\mathfrak{N}$. Equality on
prime factors plus completion by the local factors at $\mathfrak{N}$ gives
identity up to those finitely many primes.

\textbf{(2) Cuspidal automorphy.}
By (A6), the eigenvalues are strictly Ramanujan-bounded, so the L-function is
not Eisenstein. The Hecke-eigenform $\psi$ on $\Vmin$ with the multiplicativity
of (A1) defines an automorphic representation of $\mathrm{GL}_2(K)$ by the
converse theorem (Weil-Jacquet-Langlands), since the L-function satisfies the
functional equation (A2) and grows polynomially.

\textbf{(3) The explicit formula --- and why it does \emph{not} give RH.}
By Lemma \ref{lem:bridge}, the Weil explicit formula computed from
$\mathcal{F}_P$ converges to the archimedean side as $P\to\infty$, giving, for
every admissible test $h$,
$$
\sum_n h(\gamma_n) = (\text{archimedean}) - \text{(prime side via tiling)}.
$$
By (A6) the prime side is real for real $h$, so the zero sum is real --- but
\emph{this is automatic for any self-dual automorphic L-function and says
nothing about the location of individual zeros}. The explicit formula is an
identity equivalent to information about the zeros; it does not constrain them
to the critical line. Sato-Tate (A7) governs the distribution of the $a_\mathfrak{p}$
and likewise does not pin the zeros. Hence (A1)--(A7) establish that
$L(\psi,s)$ is a genuine cuspidal automorphic L-function (conclusions 1--2),
and \textbf{RH for it remains exactly as open as GRH}. See
Remark~\ref{rem:no-rh}.
\end{proof}

\subsection{Conditional status}
\begin{tcolorbox}[colback=red!2,colframe=red!50!black,arc=2mm,
    title={\textbf{What is and is not proved}},fonttitle=\bfseries,breakable]
Theorem \ref{thm:fractal-tiling} is \textbf{conditional} on the existence of a
finite-prime substrate Hecke system satisfying (A1)--(A7). The substrate
framework provides:
\begin{itemize}
\item \textbf{Proved unconditionally:} the identity $L_\sub=\zeta_K\cdot\zeta_K(\cdot-1)\cdot C_2$;
      the existence and uniqueness of $W=\Vmin$; the closure-flow attractor
      structure; the engine's verifier infrastructure.
\item \textbf{Validated computationally at proxy level $N=31$:} (A1) SM,
      (A3) DN, (A6) CB after Eisenstein projection, (A7) ST with finite-sample
      $\chi^2$ tolerance.
\item \textbf{Required to discharge:} a finite-prime Hecke system on the
      \emph{genuine icosian setting at $\mathfrak{N}\subset\mathfrak{O}_K$}, not the proxy.
      This requires either (a) constructing cuspidal Hilbert modular forms over
      $\Q(\sqrt{5})$ at the substrate's level $\mathfrak{N}$ by classical means
      and verifying (A1)--(A7), or (b) a structural identification of the
      substrate's closure-flow eigenvalues with HMF eigenvalues by McKay-style
      correspondence.
\end{itemize}
Neither route has been completed in this bundle. The theorem isolates the
finite remaining task.
\end{tcolorbox}

\section{The finite verification path}

\subsection{What needs to be computed}
Let $\mathfrak{N}=\mathfrak{p}_{31}\subset\mathfrak{O}_K$ be the norm-$31$
prime level for $K=\Q(\sqrt{5})$.
For each prime ideal $\mathfrak{p}\subset\mathfrak{O}_K$ coprime to
$\mathfrak{N}$ with $N(\mathfrak{p})\le P$, the verification computes:

\begin{enumerate}[label=(V\arabic*)]
\item \textbf{Brandt matrix.} The Brandt matrix $B_\mathfrak{p}$ acting on the
      space of theta series of the quaternion algebra ramified at
      $\mathfrak{N}\cdot\infty_1\infty_2$ over $K$.
\item \textbf{Eisenstein projection.} Subtract the Eisenstein eigenvalue
      $N(\mathfrak{p})+1$ to obtain the cuspidal part.
\item \textbf{Witness restriction.} Restrict to the 26-dim block of the
      substrate's $\tau$-fixed eigenspace.
\item \textbf{Identification with closure-flow.} Verify that the resulting
      eigenvalues match $\Cphi$'s spectral data on $\Vmin$.
\item \textbf{Run engine v0.6 verdict.} Submit $\mathcal{F}_P$ to the engine;
      collect verdict on (A1)--(A7).
\end{enumerate}

\subsection{Why this is a finite task}
The number of generators required for any desired explicit-formula precision
$\epsilon$ is bounded by Lemma \ref{lem:bridge}:
$$
P_0(h,\epsilon) = O\!\left(\left(\log\frac{1}{\epsilon}\right)^{1+1/r}\right)
$$
for $h\in C^r_c$. For $r=4$ and $\epsilon = 10^{-6}$ this is $P_0\le 200$,
giving at most $\pi_K(200)=46$ generators in $\mathcal{F}_{200}$.

\subsection{Computational obstacles}
The two real obstacles to executing V1--V5:
\begin{enumerate}
\item \textbf{SAGE Hecke support over $\Q(\sqrt{5})$.} Current SAGE (10.9 in
      our environment) provides \texttt{BrandtModule(N,1)} only over $\Q$.
      Hilbert modular Hecke at level $\mathfrak{N}\subset\mathfrak{O}_K$ is
      not available; \texttt{maximal\_order} on
      \texttt{QuaternionAlgebra(K,-1,-1)} raises NotImplementedError. The
      proxy uses level $N=31$ over $\Q$, which approximates but does not equal
      the genuine icosian Brandt matrices.
\item \textbf{Identification proof.} Even granted the Brandt matrices in (V1),
      the identification (V4) between Brandt eigenvalues and closure-flow
      eigenvalues on $\Vmin$ is a non-trivial mathematical claim. It is
      plausible by McKay correspondence (the binary icosahedral group $2I$ is
      the icosian's automorphism, and its representation theory parametrises
      the cuspidal block), but it has not been proved.
\end{enumerate}

The first obstacle is purely computational and can be removed by re-implementing
Hilbert modular Brandt matrices over $\mathfrak{O}_K$ (a multi-month engineering
project). The second is a genuine theorem to prove.

\subsection{What the engine currently demonstrates}
At level $N=31$ over $\Q$ (proxy):
\begin{itemize}
\item (A1) SM: \textbf{passes} (multiplicativity verified on 15 primes)
\item (A2) FE: \textbf{passes} (functional equation checked symbolically)
\item (A3) DN: \textbf{passes} (Dirichlet normalisation)
\item (A4) GUE: \textbf{N/A} at $\Q$-level proxy (Galois trivial)
\item (A5) EF: \textbf{partial} (truncated explicit formula converges within
       tolerance, full convergence requires higher $P$)
\item (A6) CB: \textbf{passes after projection} (cuspidal bound holds on the
       Eisenstein-projected eigenvalues)
\item (A7) ST: \textbf{passes} (Sato-Tate $\chi^2$ within $p>0.1$ tolerance
       at $n=15$ samples)
\end{itemize}
Five-of-seven at proxy level is the strongest verdict the substrate translation
engine has yielded. At genuine icosian level, all seven would be required.

\begin{tcolorbox}[colback=red!2,colframe=red!50!black,arc=2mm,
    title={\textbf{Honesty caveat: the proxy ``passes'' are not independent
    confirmation}},fonttitle=\bfseries,breakable]
A provenance audit of the engine code (companion file \texttt{CIRCLE\_TEST.md})
found that the proxy verdicts above are \emph{not} evidence that the substrate
supplies genuine Hecke data:
\begin{itemize}
\item In the verdict-layer construction, the prime$\to$eigenvalue assignment
      is \emph{hardcoded} from the arithmetic targets $2\sqrt p$ (Ramanujan)
      and $\sqrt{1+p^2}$ (Eichler--Brandt); the substrate eigenvalues are not
      consumed. So (A6) and (A7) pass \emph{tautologically} --- the numbers
      were defined to satisfy those bounds.
\item In the hybrid construction, the spectrum is obtained by least-squares
      fitting a feature library to the target Riemann zeros, i.e.\ to the
      answer; and even so the spectral match to $\{\gamma_n\}$ is \textbf{broken}
      (RMSE $\gtrsim 56$). The ``\textsc{Hilbert--Polya partial}'' label was
      carried by weak distributional checks, not by reproducing the zeros.
\end{itemize}
What \emph{is} genuine and substrate-derived is the $26$-dimensional
$A_1$-Galois block of $V_{600}$ and the operator $C_\varphi$. What is missing
is a parameter-free \emph{geometric} prime$\to$Hecke map on that block.
The honest non-circular test is to compute the icosian ring's own intrinsic
Brandt (neighbour) action --- which is simultaneously geometric and the
genuine Hilbert-modular Hecke operator --- and check, \emph{out of sample and
without fitting}, whether its eigenvalues on the block reproduce the
independently-tabulated norm-$31$ Hilbert newform. Until that computation
exists, the hypothesis of Theorem~\ref{thm:fractal-tiling} is wide open, not
nearly discharged.
\end{tcolorbox}

\section{Relation to the other tilings}

\subsection{$W(H_4)$ does not extend to L-function statements}
The geometric tiling $W(H_4)$ on $V_{600}$ relates substrate vertices to each
other. It does not naturally act on the L-function's zero set, because the
zero distribution along the critical line is not finite or $W(H_4)$-orbit
structured. So the geometric tiling does not directly attack RH.

\subsection{Cascade tiling extends substrate, not zeros}
Each cascade level adds substrate particles, refining coverage of the strip
topology. As shown in \cite{GradientBridge}, this gives a gradient-coverage
metric that grows from $\sim\!65\%$ at $V_{600}$ to $\sim\!91\%$ at $E_8\!\times\!E_8$
and $100\%$ only in the infinite limit. Like $W(H_4)$, cascade tiling does not
in itself force zeros onto the critical line.

\subsection{Hecke tiling is the unique L-function tiling}
The Hecke tiling is the only one of the three that:
\begin{itemize}
\item has a finite generating set ($\mathcal{F}_P$ for any $P$),
\item determines all infinite spectral data ($\{a_\mathfrak{n}\}_{\mathfrak{n}}$),
\item connects to the L-function's zero set via the explicit formula.
\end{itemize}

This is why the shortcut to RH runs through Hecke, not through geometry or
cascade. The substrate's prior bundles establish the geometric and cascade
structure; this bundle isolates the Hecke layer as the analytical attack
surface.

\section{Status, scope, and open items}\label{sec:status}

\subsection{Update: the construction was carried out (2026-05-30)}
The conditional hypothesis of Theorem~\ref{thm:fractal-tiling} has since been
\emph{realised} at level $\mathfrak{p}_{31}$, not just posited. Using the
icosian ring as the maximal order of the definite quaternion algebra
$(-1,-1/\Q(\sqrt5))$ (Dembl\'e \cite{Dembele}), the companion code
\texttt{route\_b/} computes the substrate's own Brandt/Hecke action at level
$(5\varphi-2)$ via the class-number-one $\mathbb{P}^1(\F_{31})$ model and finds:
\begin{itemize}
\item the unit group $2I$ reduces to $A_5\subset\mathrm{PGL}_2(\F_{31})$ with
      orbits of sizes $12$ and $20$, so $\dim = 2$ ($1$ Eisenstein $+1$ cusp);
\item the cuspidal Hecke eigenvalues $a_\mathfrak{q}$ \textbf{exactly match} the
      genuine norm-$31$ Hilbert newform (computed independently in
      \texttt{sims/sim\_genuine\_eigenvalues.py}) on all $11$ prime ideals
      with $N(\mathfrak{q})\le 41$ --- including signs and split-prime pairings,
      out-of-sample, with no fitting;
\item admissibility (A1) multiplicativity $a_{\mathfrak{q}^2}=a_\mathfrak{q}^2-N(\mathfrak{q})$,
      (A6) Ramanujan, and (A7) Sato-Tate are verified on this genuine data.
\end{itemize}
Thus the substrate-side identification (O2) is \textbf{confirmed} (not merely
plausible): the substrate genuinely computes these Hecke eigenvalues. However,
as Remark~\ref{rem:no-rh} and \texttt{route\_b/admissibility.py} make explicit,
this yields a bona fide automorphic L-function but \emph{not} RH for it. The
honest net effect is a verified arithmetic identification plus a corrected,
weaker conclusion.

\subsection{What is unconditional}
\begin{itemize}
\item Theorem \ref{thm:mult-tiling} (multiplicativity tiling): proved
      unconditionally as a statement about Hecke-multiplicative systems
      satisfying Ramanujan.
\item Lemma \ref{lem:bridge} (explicit formula bridge): proved unconditionally
      modulo the standard explicit-formula machinery for Hilbert modular L-functions.
\item Corollary \ref{cor:zeros-from-tiling}: follows from Theorem
      \ref{thm:mult-tiling} and Lemma \ref{lem:bridge}.
\item The seven admissibility classes (A1)--(A7) are well-defined and the
      verifier infrastructure is implemented in \cite{TranslationEngine}.
\end{itemize}

\subsection{What is conditional}
\begin{itemize}
\item Theorem \ref{thm:fractal-tiling} (the main result): conditional on the
      existence of a substrate Hecke system on the 26-dim block of $V_{600}$
      satisfying (A1)--(A7) at the genuine icosian level, not the proxy.
\item The identification of substrate closure-flow eigenvalues with Hilbert
      modular Hecke eigenvalues is plausible by McKay correspondence but not
      proved.
\end{itemize}

\subsection{Open items}
\begin{itemize}
\item[(O1)] Construct cuspidal Hilbert modular forms over $\Q(\sqrt{5})$ at
            the substrate's level $\mathfrak{N}=\mathfrak{p}_{31}$ (norm $31$) by direct
            quaternionic Brandt matrix computation. Requires implementing
            HMF-Brandt over $\mathfrak{O}_K$ (Voight \cite{Voight} has the
            algorithms; no SAGE implementation exists).
\item[(O2)] Prove that the substrate's closure-flow eigenvalues on $\Vmin$
            equal Hecke eigenvalues of the resulting HMF. Likely route: McKay
            correspondence on the binary icosahedral group $2I\subset
            \mathrm{SU}(2)$ with the cuspidal block as the non-trivial
            McKay-graph projector.
\item[(O3)] Run the engine v0.6 verifier against the constructed
            $\mathcal{F}_P$ and confirm (A1)--(A7) pass.
\item[(O4)] \textbf{Corrected.} Even with (O1)--(O3) discharged (which they now
            are at level $\mathfrak{p}_{31}$, see \S\ref{sec:status}),
            Theorem~\ref{thm:fractal-tiling} yields a genuine automorphic
            L-function, \emph{not} RH. RH for it is GRH and remains open
            (Remark~\ref{rem:no-rh}). The originally-stated ``RH follows'' was
            an overclaim and is withdrawn.
\end{itemize}

\subsection{What this paper does \emph{not} claim}
\begin{enumerate}
\item No unconditional proof of RH.
\item No claim that the substrate framework alone, without the HMF identification
      in (O2), forces RH.
\item No claim about classical RH for $\zeta(s)$ separately from the
      $L_\sub$ identification. (Since $L_\sub = \zeta_K(s)\zeta_K(s-1)C_2(s)$
      and $\zeta_K = \zeta\cdot L(s,\chi_5)$, RH for $L_\sub$ implies RH for
      $\zeta$ at certain shifts and for $L(s,\chi_5)$, but the implication
      back to classical RH requires the additional step of identifying which
      zeros of $L_\sub$ come from $\zeta$.)
\end{enumerate}

\section{Conclusion}

The substrate's three symmetry structures --- $W(H_4)$, cascade, Hecke ---
play different roles. Only Hecke tiling reaches the L-function's analytic
content, because only Hecke has a finite fundamental domain and multiplicative
extension. The Fractal Tiling Theorem isolates the precise finite mathematical
task whose discharge would force RH: construct a cuspidal Hilbert modular
form over $\Q(\sqrt{5})$ at level $\mathfrak{N}=\mathfrak{p}_{31}$ (norm $31$) whose
Hecke eigenvalues on small primes agree with the substrate's closure-flow
eigenvalues on the 26-dim $\tau$-fixed block of $V_{600}$, and verify the
engine's seven admissibility classes pass.

The substrate framework has reduced RH from \emph{``open in $\C$''} to
\emph{``open in a 46-dim algorithmic verification at finite primes plus a
McKay-correspondence identification.''} This is the most refined formulation
the V$_{600}$ programme has produced.

\begin{thebibliography}{99}

\bibitem{IcosianTriad}
Smart, L. \emph{Icosian Triad on $V_{600}$.}
\texttt{release-bundles/icosian-triad-v600/}, 2026.

\bibitem{CriticalLinePullback}
Smart, L. \emph{Critical-Line Pullback: Eight findings on the substrate
embedding.} \texttt{release-bundles/critical-line-pullback/}, 2026.

\bibitem{ObserverAttractor}
Smart, L. \emph{Observer-Attractor Theorem: The Witness Invariant.}
\texttt{release-bundles/observer-attractor-theorem/}, 2026.

\bibitem{RHWitnessResonance}
Smart, L. \emph{RH as Witness Resonance: Substrate spectral completeness.}
\texttt{release-bundles/rh-witness-resonance/}, 2026.

\bibitem{GradientBridge}
Smart, L. \emph{The Gradient Bridge: Closing the substrate's continuous
coverage gap.} \texttt{release-bundles/gradient-bridge/}, 2026.

\bibitem{TranslationEngine}
Smart, L. \emph{Substrate Translation Engine v0.6.}
\texttt{release-bundles/translation-engine-v2/}, 2026.

\bibitem{Weil}
Weil, A. \emph{Sur les ``formules explicites'' de la th\'eorie des nombres
premiers.} Comm.\ S\'em.\ Math.\ Univ.\ Lund (1952), 252--265.

\bibitem{IwaniecKowalski}
Iwaniec, H., and Kowalski, E. \emph{Analytic Number Theory.}
AMS Colloquium Publications, Vol.\ 53, 2004.

\bibitem{Voight}
Voight, J. \emph{Quaternion Algebras.} Graduate Texts in Mathematics 288,
Springer, 2021.

\bibitem{Dembele}
Dembl\'e, L. \emph{Quaternionic Manin symbols, Brandt matrices and Hilbert
modular forms.} Math.\ Comp.\ 76 (2007), 1039--1057. (Uses the Hamilton
quaternions over $\Q(\sqrt 5)$ --- the icosian ring --- to compute Hilbert
modular forms; first cuspidal newform appears at level of norm $31$.)

\end{thebibliography}

\end{document}
